\documentclass[11pt]{article}
\usepackage[a4paper,margin=27mm]{geometry}
\usepackage{amsmath,amssymb,mathtools}
\usepackage{fontspec,unicode-math}
\setmainfont{Latin Modern Roman}
\setmathfont{Latin Modern Math}
\usepackage{xurl}
\usepackage[hidelinks]{hyperref}
\newcommand{\sref}[2]{\hyperlink{source:#1:#2}{[S#2]}}
\newcommand{\cataloguescope}[1]{\paragraph{Mathematical question.}#1}
\hypersetup{pdftitle={Problem statement}}
\begin{document}
\section{Wikipedia number-theory item 173: For any given integer b which is not a perfect power and not of the form −4k\textasciicircum{}\{4\} for integer k, are…}\label{statement}

\subsection{Definitions and mathematical statement}

\paragraph{Shared notation.}$\mathbb N=\{1,2,\ldots\}$, and $\mathbb Z,\mathbb Q,\mathbb R,\mathbb C$ denote the integers, rationals, reals and complex numbers. Any other conventions are specified locally.


Fix an integer base $b$ with $|b|\ge2$. For $n\in\mathbb N$ define
\[R_n(b)=1+b+\cdots+b^{n-1}=\frac{b^n-1}{b-1}.\]
A repunit prime to base $b$ is a value $R_n(b)$ which is a positive prime. Here a perfect power is an integer $m^r$ with $m\in\mathbb Z$ and integer $r>1$; negative bases are included in this convention.
\paragraph{Statement recorded in the supplied source.}
For every integer $b$ with $|b|\ge2$ which is neither a perfect power nor of the form $-4k^4$ for an integer $k$, is the set
\[\{R_n(b):n\in\mathbb N,\ R_n(b)\text{ is a positive prime}\}\]
infinite? \sref{statement}{2}
\subsection{Short English statement}
For every admissible positive or negative base, do its repunits take infinitely many positive prime values?
\subsection{Sources}
\begin{description}
\item[\hypertarget{source:statement:1}{S1}] ULAM.AI and the UnsolvedMath Contributors, UnsolvedMath: A Curated Collection of Open Mathematics Problems, catalogue entry. CC BY 4.0. \url{https://huggingface.co/datasets/ulamai/UnsolvedMath}.
\item[\hypertarget{source:statement:2}{S2}] Wikipedia contributors, \emph{List of unsolved problems in mathematics}, revision 1366636767 (2026) Number theory, Wikipedia number-theory item 173: For any given integer b which is not a perfect power and not of the form −4k\textasciicircum{}\{4\} for integer k, are…. \url{https://en.wikipedia.org/w/index.php?title=List_of_unsolved_problems_in_mathematics&oldid=1366636767#Number_theory}.
\item[\hypertarget{source:statement:3}{S3}] Wikipedia contributors, \emph{Repunit}, revision 1363388346 (2026) Definition; Algebra factorization of generalized repunit numbers, infinitude conjecture paragraph. \url{https://en.wikipedia.org/w/index.php?title=Repunit&oldid=1363388346}.
\end{description}
\label{end:statement}
\end{document}
